\section{Part I: Tensor Network Representations}
\label{sec:part1}
\lecture{28}{09}{2026}

% Mirrors slides.qmd "Part I" outline. One subsection per outline item.

Everything in this course rests on a single bet. A field sampled on a fine grid
is an enormous vector, but it is not an \emph{arbitrary} enormous vector: it is
smooth, it is multiscale, and the coupling between well-separated scales is
weak. If that is true, the vector lies on a manifold of far lower dimension than
the space it sits in, and we should compute on the manifold rather than in the
ambient space.

Part~\partref{1} builds the machinery needed to state that bet precisely: a
diagrammatic language (\cref{subsec:penrose-notation}), a short account of where
the ideas came from (\cref{subsec:literature-survey}), the scaling wall they are
meant to break (\cref{subsec:numerical-methods-intro}), the ansatz itself
(\cref{subsec:mps}), the binary re-indexing that makes it multiscale
(\cref{subsec:qtt-functions-operators}), and --- by the end --- a working PDE
solver whose cost is logarithmic in the number of grid points
(\cref{subsec:heat-explicit-euler}).

\subsection{Penrose Diagram Notation}
\label{subsec:penrose-notation}
% Keywords: tensor as a node, legs = indices, contraction = joined legs,
% scalars/vectors/matrices/tensors as diagrams, index summation convention,
% trace, identity/copy tensor, diagrammatic vs. index notation.

Index expressions with five or six summations are correct and unreadable.
Diagrams are the working notation of the field for a simple reason: they carry
exactly the same information, and they make the \emph{structure} of a
computation something you can see.
\margintext{Diagrammatic tensor notation is due to Penrose (1971); it entered
numerical practice through the DMRG community in the 1990s and is now standard
in both communities described in \cref{subsec:literature-survey}.}

\begin{itemize}
  \item A tensor is a \textbf{node}; every index is a \textbf{leg}. The rank of
    a tensor is simply its number of legs.
  \item A scalar has no legs, a vector $v_i$ one, a matrix $M_{ij}$ two, a
    general tensor $T_{ijk}$ three, and so on.
  \item Nothing is lost in the translation --- the diagram carries precisely
    what the index expression carries.
\end{itemize}

There are only two moves, and every construction in this course is assembled
from them: \emph{place nodes} and \emph{join legs}.

\begin{itemize}
  \item Joining two legs means \textbf{summing over that shared index}:
    \begin{equation*}
      \sum_j M_{ij} N_{jk}
      \quad\longleftrightarrow\quad
      \text{two nodes, one bond.}
    \end{equation*}
  \item \textbf{Dangling legs} are the free indices of the result. Counting them
    tells you what object you have built.
  \item A \textbf{closed loop} has no free legs left, so the diagram evaluates to
    a number: it is a trace.
\end{itemize}

A small vocabulary of recurring nodes covers most diagrams you will meet:

\begin{itemize}
  \item \textbf{Identity} --- a bare line.
  \item \textbf{Copy (or $\delta$) tensor} --- a branching node, one index shared
    by three legs.
  \item \textbf{Isometry} --- drawn shaded; contracted with its conjugate it
    gives the identity.
\end{itemize}

\begin{notation}
  Throughout these notes physical legs point \textbf{down} and virtual (bond)
  legs run \textbf{horizontally}. Conventions differ between papers; the
  diagrams only mean something once the convention is fixed.
\end{notation}

\begin{remark}
  The real payoff is not compactness. It is that the \emph{order} of contraction
  is visible in the diagram --- and contraction order is what determines cost.
  Two evaluation orders of the same diagram can differ by many orders of
  magnitude in runtime while producing identical numbers.
\end{remark}

\subsection{Literature Survey}
\label{subsec:literature-survey}
% Keywords: DMRG (White 1992), matrix product states, tensor train
% decomposition (Oseledets), quantics/quantized tensor trains (QTT),
% tensor networks for PDEs, tensor networks for CFD, review articles.

Before the technical material, a short orientation --- partly because the
subject has an unusual history. The same ansatz was invented twice, in two
communities that did not talk to each other, and \emph{both} vocabularies remain
in active use. You will meet both in this course and in the papers we draw on.

\begin{center}
  \begin{tabular}{ll}
    \toprule
    \textsc{Quantum many-body} & \textsc{Numerical analysis} \\
    \midrule
    Matrix product state (MPS)   & Tensor train (TT) \\
    Bond dimension $\chi$        & TT-rank \\
    Entanglement across a cut    & Singular-value decay \\
    DMRG sweep                   & Alternating least squares \\
    \bottomrule
  \end{tabular}
\end{center}

\noindent
The two columns denote the same objects and the same algorithms.

\paragraph{Where it came from.}
\margintext{Note the seventeen-year gap between White (1992) and Oseledets
(2011): the numerical-analysis rediscovery was independent, and only later were
the dictionaries matched up.}
\begin{itemize}
  \item \textbf{DMRG} --- variational optimization over matrix product states,
    born in condensed-matter physics~\cite{white1992}.
  \item \textbf{MPS in modern language} --- the standard reference for canonical
    forms and sweeping~\cite{schollwoeck2011}.
  \item \textbf{Tensor-train decomposition} --- the same ansatz rederived for
    numerical linear algebra, with its own error analysis~\cite{oseledets2011}.
  \item \textbf{Quantics} --- binary indexing gives $O(d \log N)$ approximation
    of $N$-$d$ tensors~\cite{khoromskij2011}.
  \item \textbf{TT-cross} --- build a tensor train from samples, without ever
    forming the full array~\cite{oseledets2010ttcross}.
\end{itemize}

\paragraph{Tensor networks meet fluids.}
\begin{itemize}
  \item \textbf{Turbulence has exploitable structure.} Correlations between
    widely separated scales are weak, so the bond dimension stays
    modest~\cite{gourianov2022}.
  \item \textbf{Kinetic plasmas.} Vlasov--Poisson solved in MPS form with large
    compression~\cite{ye2022vlasov}.
  \item \textbf{Reduced-order models.} Wall-bounded flows represented at a few
    percent of the DNS variable count~\cite{kiffner2023rom}.
  \item \textbf{A full stack.} Incompressible flow around arbitrary immersed
    geometry, end to end~\cite{peddinti2024cfd}.
\end{itemize}

\paragraph{Where the field is now.}
\margintext{Our own route through the subject: the 2024 framework fixed the
encoding and the projection loop; the 2026 turbulence study asked how far the
ranks stretch in 3D; the 2025 space--time paper removed the time loop
altogether. Parts~\partref{3} and~\partref{4} follow that order.}
\begin{itemize}
  \item \textbf{3D turbulence} --- compression, simulation and synthesis at
    $\mathrm{Re}_\lambda = 315$ on a $1024^3$ grid~\cite{pisoni2026turbulence}.
  \item \textbf{Compressible flow} --- a 2D Euler solver, finalist in the
    Airbus--BMW Quantum Computing Challenge 2024~\cite{peddinti2025compressible}.
  \item \textbf{Space--time and data-driven} --- space and time in a single MPS,
    together with MPS-DMD~\cite{peddinti2025spacetime}.
\end{itemize}

Parts~\partref{3} and~\partref{4} are built directly on these last
three.

\subsection{Numerical Methods: A Brief Introduction}
\label{subsec:numerical-methods-intro}
% Keywords: curse of dimensionality, grid-based discretization, finite
% differences, why classical solvers struggle at high resolution/dimension,
% compressed representations, low-rank structure.

The motivation for all of this is a scaling wall, so it is worth stating the
wall carefully before claiming to climb it.

\paragraph{The curse of dimensionality.}
\begin{itemize}
  \item Discretize a field on a mesh with $N = 2^n$ points per dimension: in $d$
    dimensions that is $2^{nd}$ unknowns.
  \item A $1024^3$ mesh is already $\sim\!10^9$ unknowns --- \emph{per field,
    per timestep}.
  \item In CFD the cost driver is resolution, and three-dimensional direct
    numerical simulation needs $\sim\!\mathrm{Re}^{9/4}$ degrees of freedom
    (\cref{subsec:navier-stokes}).
  \item Exascale hardware does not close this gap. It moves it by a couple of
    orders of magnitude.
\end{itemize}

\paragraph{Why structure is the only way out.}
\margintext{Every classical accelerator --- adaptive meshes, multigrid, spectral
methods, POD and reduced-order models --- is an exploitation of structure. The
tensor-network ansatz is another one, not a different kind of thing.}
\begin{itemize}
  \item The standard toolkit --- finite differences, volumes or elements,
    explicit or implicit stepping, a linear solve per step --- does not change
    the scaling. It changes the constant.
  \item \textbf{The tensor-network bet:} the exponentially long vector is not an
    arbitrary point in $\mathbb{R}^{2^n}$. A smooth, multiscale field lives on a
    low-rank manifold described by $O(n\chi^2)$ numbers.
  \item \textbf{The price:} exact linear algebra is replaced by approximate
    arithmetic, with a controlled truncation at every single step.
\end{itemize}

That trade --- exponential memory for a truncation error you must monitor --- is
the recurring theme of the entire course.

\subsection{Matrix Product States}
\label{subsec:mps}
% Keywords: MPS/tensor train ansatz, bond dimension, canonical form
% (left/right/mixed), SVD-based compression, singular value decay,
% area law of entanglement, exact vs. approximate ranks, orthogonality center.

\begin{definition}[Matrix product state / tensor train]
  A vector with $2^n$ components is written as a chain of rank-3 tensors,
  \begin{equation}
    v_{i_1 i_2 \cdots i_n}
      = A^{(i_1)}_1 A^{(i_2)}_2 \cdots A^{(i_n)}_n ,
    \label{eq:mps}
  \end{equation}
  where each $A^{(i_k)}_k$ is a matrix of size at most $\chi \times \chi$. The
  \textbf{bond dimension} $\chi$ is the size of the virtual legs.
\end{definition}

The parameter count of~\eqref{eq:mps} is at most $2n\chi^2$ rather than $2^n$.
That is exponential compression --- \emph{provided $\chi$ stays bounded as $n$
grows}. Essentially everything that follows is about when that proviso holds.

\paragraph{Why an MPS can work at all.}
\begin{itemize}
  \item Cut the chain at bond $k$. The bond dimension needed there is the rank of
    $v$ reshaped into a matrix across that cut.
  \item That rank is governed by \textbf{singular-value decay} --- which
    physicists read as the entanglement between the two halves.
  \item For ground states of local Hamiltonians an \textbf{area law} bounds the
    entanglement independently of system size, hence $\chi$ independent of $n$.
  \item Smooth, multiscale functions behave the same way. Rough, noisy or
    randomly structured data do not.
\end{itemize}

\paragraph{Building an MPS: TT-SVD.}
\begin{itemize}
  \item Reshape the vector, take an SVD, split off one tensor, and repeat along
    the chain.
  \item Truncating small singular values at each cut is
    \textbf{quasi-optimal}: within a factor $\sqrt{n}$ of the best rank-$\chi$
    approximation~\cite{oseledets2011}.
  \item Exact and fully controlled --- but it needs the dense $2^n$ vector in
    memory, so it is practical only for modest $n$.
\end{itemize}

\Cref{subsec:finding-qtt} returns to this with methods that build an MPS
\emph{without} ever forming the full array; that is what makes the whole
programme viable.

\paragraph{Canonical forms and rounding.}
\margintext{Gauge freedom is not a nuisance to be tolerated; fixing it is what
makes truncation optimal and inner products cheap. In practice one carries the
orthogonality centre along with the state.}
\begin{itemize}
  \item The factorization~\eqref{eq:mps} is not unique: insert $G G^{-1}$ on any
    bond and nothing changes. Fixing this \textbf{gauge freedom} gives the
    canonical forms.
  \item Sweeping QR from the left and from the right produces \textbf{left,
    right} and \textbf{mixed} canonical forms, the last carrying an
    orthogonality centre.
  \item \textbf{Rounding is the workhorse of the whole subject.} Every operation
    inflates ranks --- a product of two MPSs has bond dimension $\chi_1\chi_2$
    --- so one recompresses back to $\chi$ once per step, every step.
\end{itemize}

\subsection{Quantics Tensor Trains: Functions and Operators}
\label{subsec:qtt-functions-operators}
% Keywords: binary/quantics encoding of a grid, QTT representation of
% functions, exponential compression, QTT rank vs. function smoothness,
% representing differential/difference operators in QTT form, QTT of
% derivative/Laplacian operators.

So far an MPS is a way of factorizing a long vector. The quantics trick is what
turns it into a statement about \emph{physics}, by making each tensor in the
chain correspond to a length scale.

\begin{definition}[Quantics encoding]
  Index a grid point by its binary digits,
  \begin{equation}
    x = \sum_{k=1}^{n} i_k 2^{-k}, \qquad i_k \in \{0,1\}.
    \label{eq:quantics}
  \end{equation}
  A length-$2^n$ vector is thereby reshaped into an $n$-leg tensor, which is then
  written as an MPS. The result is a \textbf{quantics tensor train} (QTT).
\end{definition}

\begin{itemize}
  \item Each tensor now carries \textbf{one length scale}: leg~1 the coarsest,
    leg~$n$ the finest.
  \item The chain becomes a multiscale, wavelet-like object, and the bond between
    legs $k$ and $k+1$ measures the coupling \emph{between scales}.
\end{itemize}

\paragraph{Which functions are low-rank?}
\begin{itemize}
  \item Exact and tiny: $e^{ax}$, $\sin$ and $\cos$ all have $\chi = 2$.
  \item Polynomials of degree $p$ have $\chi = p+1$.
  \item Sums and products of these stay low-rank: ranks add and multiply, and
    then you round.
  \item Doubling the resolution costs $O(n)$ additional parameters rather than
    $O(2^n)$.\margintext{The rule of thumb worth memorizing: rank tracks
    smoothness and scale separation, \emph{not} resolution.}
\end{itemize}

\paragraph{Bit ordering in two and three dimensions.}
With more than one coordinate there is a genuine choice to make:

\begin{itemize}
  \item \textbf{Stacked} --- $x_1 x_2 \cdots x_n\, y_1 y_2 \cdots y_n$: all $x$
    scales, then all $y$ scales.
  \item \textbf{Interleaved} --- $x_1 y_1 x_2 y_2 \cdots$: the coarsest scale of
    both coordinates, then the next, and so on.
\end{itemize}

Interleaving keeps each \emph{scale} local along the chain, which matches the
physics of a turbulent cascade. This is not a cosmetic choice: it decides
whether non-Gaussian turbulence statistics survive
compression~\cite{pisoni2026turbulence}. We return to the evidence in
\cref{subsec:expressivity}.

\paragraph{Operators as MPOs.}
\begin{definition}[Matrix product operator]
  The operator analogue of an MPS: a chain in which each tensor carries two
  physical legs (one in, one out) in addition to its two virtual bonds.
\end{definition}

\begin{itemize}
  \item Finite-difference shift, first-derivative and Laplacian operators have
    \textbf{exact and tiny} MPO ranks, $\chi_{\mathrm{MPO}} \approx 2\text{--}4$.
  \item Applying an MPO to an MPS costs
    $O(n \chi^2 \chi_{\mathrm{MPO}}^2)$, plus a rounding step.
  \item The dense $2^n \times 2^n$ matrix is never formed --- it never exists in
    memory at any point.
\end{itemize}

\subsection{Solving the Heat Equation (Explicit Euler Method)}
\label{subsec:heat-explicit-euler}
% Keywords: 1D/2D heat equation, explicit time-stepping, stability
% (CFL-type) condition, QTT representation of the state, applying the
% discrete Laplacian in QTT form, rounding/recompression after each step,
% numerical example.

We now have enough to solve a PDE. The heat equation is deliberately the easiest
possible choice --- diffusion smooths, so the ranks ought to behave --- which
makes it the right place to see the machinery work and to isolate what can go
wrong.

\paragraph{The setup.}
Take $\partial_t u = \alpha\, \partial_x^2 u$ on $[0,1]$, discretized with
$N = 2^n$ points. Three ingredients suffice, and there is nothing else:

\begin{itemize}
  \item the initial condition $u^0$ as an MPS;
  \item the discrete Laplacian $L$ as an MPO;
  \item one scalar, $\alpha \Delta t / \Delta x^2$.
\end{itemize}

\paragraph{One timestep.}
Explicit Euler is a single line of tensor arithmetic,
\begin{equation}
  u^{k+1} = \left(\mathbb{1} + \alpha \Delta t\, L\right) u^{k},
  \label{eq:heat-explicit}
\end{equation}
consisting of exactly two operations: an \textbf{MPO--MPS product}, and then
\textbf{rounding} to a tolerance $\varepsilon$ (or to a fixed $\chi$).

\begin{itemize}
  \item Cost is $O(n\chi^3)$ per step with $O(n\chi^2)$ memory --- both
    \textbf{logarithmic in the mesh size}.
  \item Skipping the rounding is the classic first mistake: ranks double at every
    step, and within twenty steps you are back to storing $2^n$ numbers.
\end{itemize}

\paragraph{What to watch.}
\margintext{Two error sources, kept apart: discretization
$O(\Delta t, \Delta x^2)$ and truncation. Monitor both --- and always know which
one currently dominates.}
\begin{itemize}
  \item \textbf{Compression does not rescue stability.} The condition is still
    $\alpha \Delta t / \Delta x^2 \le 1/2$, and with $\Delta x = 2^{-n}$ that is
    punishing. This is precisely the motivation for
    \cref{subsec:heat-implicit-euler}.
  \item \textbf{Watch $\chi(t)$.} Diffusion smooths the solution, so the bond
    dimension should stay flat or fall. If it climbs, something is wrong with
    the encoding or with the rounding tolerance.
\end{itemize}

\begin{exercise}
  Run the companion notebook: explicit Euler at $n = 20$, i.e.\ $10^6$ grid
  points, on a laptop. Then remove the rounding step and observe how many
  timesteps you get before memory runs out.
\end{exercise}
